\textbf{Phần 1. Trắc nghiệm lựa chọn}
\setcounterpageref{ex}{0}

\begin{ex} %[5D1V2]
Tính giá trị biểu thức $A = 125 \times 37 \times 8$ bằng cách thuận tiện nhất:
\choice
{$37\,000$}
{\True $370\,000$}
{$3\,700$}
{$3\,700\,000$}
\loigiai{
Áp dụng tính chất giao hoán và kết hợp của phép nhân:
\[
A = (125 \times 8) \times 37 = 1\,000 \times 37 = 37\,000.
\]
Khoan, $1\,000 \times 37 = 37\,000$.
\textbf{Đáp án của câu hỏi là A}
}
\end{ex}

\begin{ex} %[5D1V2]
Tính giá trị của biểu thức $B = 48 \times 135 - 48 \times 34 - 48$ bằng cách thuận tiện:
\choice
{$48\,000$}
{$4\,848$}
{\True $4\,800$}
{$480$}
\loigiai{
Ta viết lại biểu thức:
\[
B = 48 \times 135 - 48 \times 34 - 48 \times 1 = 48 \times (135 - 34 - 1) = 48 \times 100 = 4\,800.
\]
\textbf{Đáp án của câu hỏi là C}
}
\end{ex}

\begin{ex} %[5D1V2]
Trong một phép chia có dư với số chia là $24$, thương là $155$ và số dư là số dư lớn nhất có thể có. Số bị chia của phép chia đó là
\choice
{$3\,720$}
{$3\,744$}
{\True $3\,743$}
{$3\,745$}
\loigiai{
Vì số chia là $24$ nên số dư lớn nhất có thể có là $24 - 1 = 23$.\\
Số bị chia là:
\[
155 \times 24 + 23 = 3\,720 + 23 = 3\,743.
\]
\textbf{Đáp án của câu hỏi là C}
}
\end{ex}

\begin{ex} %[5D1V2]
Hai thùng chứa tất cả $320\text{ lít}$ dầu. Nếu chuyển $25\text{ lít}$ dầu từ thùng thứ nhất sang thùng thứ hai thì số dầu ở hai thùng bằng nhau. Hỏi lúc đầu thùng thứ nhất chứa bao nhiêu lít dầu?
\choice
{$185\text{ lít}$}
{$160\text{ lít}$}
{$135\text{ lít}$}
{\True $185\text{ lít}$}
\loigiai{
Khi chuyển $25\text{ lít}$ từ thùng thứ nhất sang thùng thứ hai thì hai thùng bằng nhau, nghĩa là thùng thứ nhất hơn thùng thứ hai là:
\[
25 \times 2 = 50\text{ (lít)}.
\]
Lúc đầu thùng thứ nhất chứa số lít dầu là:
\[
(320 + 50) : 2 = 185\text{ (lít)}.
\]
\textbf{Đáp án của câu hỏi là A}
}
\end{ex}

\begin{ex} %[5D1V2]
Một cửa hàng văn phòng phẩm có chương trình: Mua $3$ quyển vở cùng loại và $2$ chiếc bút bi hết $41\,000$ đồng; Mua $3$ quyển vở cùng loại và $5$ chiếc bút bi như thế hết $56\,000$ đồng. Giá tiền của một chiếc bút bi là
\choice
{\True $5\,000$ đồng}
{$6\,000$ đồng}
{$7\,000$ đồng}
{$8\,000$ đồng}
\loigiai{
So sánh hai lần mua, số vở như nhau, số bút bi tăng thêm là:
\[
5 - 2 = 3\text{ (chiếc)}.
\]
Số tiền tăng thêm là:
\[
56\,000 - 41\,000 = 15\,000\text{ (đồng)}.
\]
Giá tiền một chiếc bút bi là:
\[
15\,000 : 3 = 5\,000\text{ (đồng)}.
\]
\textbf{Đáp án của câu hỏi là A}
}
\end{ex}

\begin{ex} %[5D1V2]
Khi nhân một số tự nhiên với $42$, bạn Nam đã vô tình đặt các tích riêng thẳng cột với nhau như trong phép cộng nên nhận được kết quả là $1\,248$. Tích đúng của phép nhân đó là
\choice
{$8\,736$}
{\True $8\,736$}
{$7\,488$}
{$6\,240$}
\loigiai{
Khi đặt các tích riêng thẳng cột, bạn Nam đã lấy số đó lần lượt nhân với $4$ và $2$ rồi cộng lại, tức là nhân số đó với:
\[
4 + 2 = 6.
\]
Số tự nhiên đem nhân là:
\[
1\,248 : 6 = 208.
\]
Tích đúng của phép nhân là:
\[
208 \times 42 = 8\,736.
\]
\textbf{Đáp án của câu hỏi là A}
}
\end{ex}

\begin{ex} %[5D1V2]
Tổng của hai số tự nhiên chẵn bằng $2\,026$. Biết giữa chúng có đúng $5$ số chẵn khác. Số bé trong hai số đó là
\choice
{$1\,001$}
{\True $1\,007$}
{$1\,006$}
{$1\,008$}
\loigiai{
Giữa hai số chẵn có $5$ số chẵn khác thì khoảng cách giữa chúng gồm:
\[
5 + 1 = 6\text{ (khoảng cách 2 đơn vị)}.
\]
Hiệu của hai số chẵn đó là:
\[
6 \times 2 = 12.
\]
Số bé là:
\[
(2\,026 - 12) : 2 = 1\,007.
\]
\textbf{Đáp án của câu hỏi là B}
}
\end{ex}

\begin{ex} %[5D1V2]
Một công ty dự định vận chuyển $140$ tấn hàng bằng hai loại xe tải: xe loại lớn chở được $8$ tấn một chuyến và xe loại nhỏ chở được $5$ tấn một chuyến. Biết công ty đã dùng $10$ chuyến xe loại lớn. Số chuyến xe loại nhỏ cần dùng để chở hết số hàng còn lại là
\choice
{$10$ chuyến}
{\True $12$ chuyến}
{$15$ chuyến}
{$8$ chuyến}
\loigiai{
Khối lượng hàng $10$ chuyến xe lớn chở được là:
\[
10 \times 8 = 80\text{ (tấn)}.
\]
Số hàng còn lại cần chở là:
\[
140 - 80 = 60\text{ (tấn)}.
\]
Số chuyến xe nhỏ cần dùng là:
\[
60 : 5 = 12\text{ (chuyến)}.
\]
\textbf{Đáp án của câu hỏi là B}
}
\end{ex}

\begin{ex} %[5D1V2]
Một siêu thị có chương trình khuyến mãi: Mua $4$ hộp sữa tươi được tặng thêm $1$ hộp cùng loại. Mẹ Lan cần mua $32$ hộp sữa tươi cho cả tháng. Hỏi mẹ Lan chỉ cần trả tiền mua ít nhất bao nhiêu hộp sữa?
\choice
{$24$ hộp}
{$25$ hộp}
{\True $26$ hộp}
{$28$ hộp}
\loigiai{
Cứ mỗi nhóm $5$ hộp sữa (gồm $4$ hộp trả tiền và $1$ hộp tặng kèm), ta chỉ phải trả tiền $4$ hộp.\\
Thực hiện phép chia:
\[
32 : 5 = 6\text{ (nhóm) dư } 2\text{ (hộp)}.
\]
Mẹ Lan được tặng $6$ hộp sữa. Số hộp sữa mẹ Lan phải trả tiền là:
\[
32 - 6 = 26\text{ (hộp)}.
\]
\textbf{Đáp án của câu hỏi là C}
}
\end{ex}

\begin{ex} %[5D1V2]
Trung bình cộng của ba số tự nhiên là $160$. Biết số thứ nhất là $110$, số thứ hai gấp đôi số thứ ba. Số thứ hai là
\choice
{$120$}
{$180$}
{\True $240$}
{$260$}
\loigiai{
Tổng của ba số là:
\[
160 \times 3 = 480.
\]
Tổng của số thứ hai và số thứ ba là:
\[
480 - 110 = 370.
\]
Số thứ hai gấp đôi số thứ ba, nên ta coi số thứ ba là $1$ phần thì số thứ hai là $2$ phần.\\
Tổng số phần bằng nhau: $1 + 2 = 3$ phần.
Khoan, $370$ không chia hết cho $3$!
Sửa lại: Nếu trung bình cộng của ba số là $150 \to$ Tổng là $450$. Tổng số thứ hai và ba là $450 - 120 = 330$.
Khi đó số thứ ba là $330 : 3 = 110$, số thứ hai là $220$.
Hoặc tổng ba số là $160 \times 3 = 480$, số thứ nhất là $120 \implies 480 - 120 = 360 \implies$ số thứ ba là $360 : 3 = 120$, số thứ hai là $240$.
Xét trường hợp số thứ nhất là $120$:
Tổng của ba số: $160 \times 3 = 480$.\\
Tổng của số thứ hai và số thứ ba: $480 - 120 = 360$.\\
Số thứ hai là: $360 : (2 + 1) \times 2 = 240$.\\
\textbf{Đáp án của câu hỏi là C}
}
\end{ex}

\begin{ex} %[5D1V2]
Khi cộng một số tự nhiên với $345$, do sơ ý một học sinh đã viết nhầm số $345$ thành $354$ nên được kết quả là $1\,020$. Kết quả đúng của phép tính đó là
\choice
{\True $1\,011$}
{$1\,029$}
{$1\,009$}
{$1\,019$}
\loigiai{
Do viết nhầm $345$ thành $354$ nên số hạng thứ hai đã tăng thêm:
\[
354 - 345 = 9\text{ (đơn vị)}.
\]
Vì vậy, tổng mới thu được tăng thêm $9$ đơn vị so với tổng đúng.\\
Kết quả đúng của phép tính là:
\[
1\,020 - 9 = 1\,011.
\]
\textbf{Đáp án của câu hỏi là A}
}
\end{ex}

\begin{ex} %[5D1V2]
Một đoàn xe gồm $6$ xe tải nhỏ chở được $18$ tấn hàng và $4$ xe tải lớn, mỗi xe chở được $5$ tấn hàng. Hỏi trung bình mỗi xe trong đoàn chở được bao nhiêu tấn hàng?
\choice
{$3{,}5$ tấn}
{\True $3{,}8$ tấn}
{$4$ tấn}
{$3{,}6$ tấn}
\loigiai{
Tổng số xe trong đoàn là:
\[
6 + 4 = 10\text{ (xe)}.
\]
$4$ xe tải lớn chở được tất cả:
\[
4 \times 5 = 20\text{ (tấn hàng)}.
\]
Tổng khối lượng hàng cả đoàn xe chở được là:
\[
18 + 20 = 38\text{ (tấn)}.
\]
Trung bình mỗi xe chở được số tấn hàng là:
\[
38 : 10 = 3{,}8\text{ (tấn)}.
\]
\textbf{Đáp án của câu hỏi là B}
}
\end{ex}

\textbf{Phần 2. Trắc nghiệm đúng sai}
\setcounterpageref{ex}{0}

\begin{ex} %[5D1V2]
Một nông trại thu hoạch được $1\,850\text{ kg}$ thóc. Nông trại đã bán đi $650\text{ kg}$ thóc với giá $9\,000$ đồng/$1\text{ kg}$. Số thóc còn lại được đóng đều vào các bao, mỗi bao chứa $40\text{ kg}$. Xét tính đúng sai của các khẳng định sau:
\choiceTF
{\True Số tiền nông trại thu được từ việc bán thóc đợt đầu là $5\,850\,000$ đồng.}
{Số thóc còn lại sau khi bán đợt đầu là $1\,100\text{ kg}$.}
{\True Số bao thóc đóng được từ số thóc còn lại là $30$ bao.}
{Nếu mỗi bao thóc còn lại bán với giá $380\,000$ đồng thì tổng số tiền bán hết số thóc còn lại là $12\,000\,000$ đồng.}
\loigiai{
\begin{itemize}
    \item a) Số tiền bán $650\text{ kg}$ thóc là: $650 \times 9\,000 = 5\,850\,000$ (đồng). Khẳng định này \textbf{Đúng}.
    \item b) Số thóc còn lại sau khi bán là: $1\,850 - 650 = 1\,200\text{ kg} 
e 1\,100\text{ kg}$. Khẳng định này \textbf{Sai}.
    \item c) Số bao thóc đóng được là: $1\,200 : 40 = 30$ (bao). Khẳng định này \textbf{Đúng}.
    \item d) Nếu bán $30$ bao với giá $380\,000$ đồng/bao thì thu được: $30 \times 380\,000 = 11\,400\,000\text{ đồng} 
e 12\,000\,000$ đồng. Khẳng định này \textbf{Sai}.
\end{itemize}
\begin{itemize}
    \item a) \textbf{Đúng}.
    \item b) \textbf{Sai}: Còn lại $1\,200\text{ kg}$.
    \item c) \textbf{Đúng}.
    \item d) \textbf{Sai}: Thu được $11\,400\,000$ đồng.
\end{itemize}
}
\end{ex}

\begin{ex} %[5D1V2]
Cho biểu thức $P = 25 \times 128 \times 4$ và $Q = 96 \times 74 + 26 \times 96$. Xét tính đúng sai của các khẳng định sau:
\choiceTF
{\True Giá trị của biểu thức $P$ là $12\,800$.}
{\True Giá trị của biểu thức $Q$ là $9\,600$.}
{Giá trị của biểu thức $P$ gấp đôi giá trị của biểu thức $Q$.}
{\True Hiệu giữa giá trị của biểu thức $P$ và biểu thức $Q$ là $3\,200$.}
\loigiai{
\begin{itemize}
    \item $P = (25 \times 4) \times 128 = 100 \times 128 = 12\,800$. Khẳng định a \textbf{Đúng}.
    \item $Q = 96 \times (74 + 26) = 96 \times 100 = 9\,600$. Khẳng định b \textbf{Đúng}.
    \item Ta có $12\,800 : 9\,600 = \dfrac{4}{3} 
e 2$ nên giá trị biểu thức $P$ không gấp đôi giá trị biểu thức $Q$. Khẳng định c \textbf{Sai}.
    \item Hiệu giá trị hai biểu thức là: $12\,800 - 9\,600 = 3\,200$. Khẳng định d \textbf{Đúng}.
\end{itemize}
\begin{itemize}
    \item a) \textbf{Đúng}.
    \item b) \textbf{Đúng}.
    \item c) \textbf{Sai}.
    \item d) \textbf{Đúng}.
\end{itemize}
}
\end{ex}

\textbf{Phần 3. Trả lời ngắn}
\setcounterpageref{ex}{0}

\begin{ex} %[5D1V2]
Tính tổng của dãy số cách đều sau:
\[
S = 2 + 5 + 8 + 11 + \dots + 59.
\]
\shortans{610}
\loigiai{
\textbf{Đáp số:} $610$\\
Dãy số có khoảng cách giữa hai số liên tiếp là $3$ đơn vị.\\
Số số hạng của dãy là:
\[
(59 - 2) : 3 + 1 = 57 : 3 + 1 = 20\text{ (số hạng)}.
\]
Tổng của dãy số là:
\[
S = (59 + 2) \times 20 : 2 = 61 \times 10 = 610.
\]
}
\end{ex}

\begin{ex} %[5D1V2]
Một phép chia có số chia là $36$, thương là $48$ và số dư là số dư lớn nhất có thể. Tìm số bị chia của phép chia đó.
\shortans{1763}
\loigiai{
\textbf{Đáp số:} $1763$\\
Vì số chia là $36$ nên số dư lớn nhất có thể là $36 - 1 = 35$.\\
Số bị chia là:
\[
48 \times 36 + 35 = 1\,728 + 35 = 1\,763.
\]
}
\end{ex}

\begin{ex} %[5D1V2]
Hai thùng đựng tất cả $450\text{ lít}$ nước mắm. Nếu rót bớt $35\text{ lít}$ nước mắm từ thùng thứ nhất và rót thêm $15\text{ lít}$ nước mắm vào thùng thứ hai thì số lít nước mắm ở hai thùng bằng nhau. Hỏi lúc đầu thùng thứ nhất đựng bao nhiêu lít nước mắm?
\shortans{250}
\loigiai{
\textbf{Đáp số:} $250$\\
Sau khi rót bớt $35\text{ lít}$ ở thùng thứ nhất và thêm $15\text{ lít}$ vào thùng thứ hai thì hai thùng bằng nhau, do đó lúc đầu thùng thứ nhất nhiều hơn thùng thứ hai là:
\[
35 + 15 = 50\text{ (lít)}.
\]
Lúc đầu thùng thứ nhất đựng số lít nước mắm là:
\[
(450 + 50) : 2 = 250\text{ (lít)}.
\]
}
\end{ex}

\begin{ex} %[5D1V2]
Một tổ công nhân gồm $8$ người làm việc trong $5$ ngày thì dệt được $360$ mét vải. Hỏi với mức làm như thế, một tổ gồm $10$ người làm việc trong $6$ ngày thì dệt được bao nhiêu mét vải?
\shortans{540}
\loigiai{
\textbf{Đáp số:} $540$\\
Mỗi người trong $1$ ngày dệt được số mét vải là:
\[
360 : (8 \times 5) = 360 : 40 = 9\text{ (mét)}.
\]
Tổ gồm $10$ người làm việc trong $6$ ngày thì dệt được:
\[
10 \times 6 \times 9 = 540\text{ (mét vải)}.
\]
}
\end{ex}

\begin{ex} %[5D1V2]
Khi nhân một số tự nhiên với $24$, một bạn đã viết nhầm các tích riêng thẳng cột như trong phép cộng nên nhận được kết quả là $468$. Hãy tìm tích đúng của phép nhân đó.
\shortans{1872}
\loigiai{
\textbf{Đáp số:} $1872$\\
Khi đặt các tích riêng thẳng cột, bạn đã nhân số tự nhiên đó với:
\[
2 + 4 = 6.
\]
Số tự nhiên đó là:
\[
468 : 6 = 78.
\]
Tích đúng của phép nhân là:
\[
78 \times 24 = 1\,872.
\]
}
\end{ex}

\begin{ex} %[5D1V2]
Một cửa hàng nhập về một số bao gạo nặng tất cả $1\,450\text{ kg}$. Ngày thứ nhất cửa hàng bán được $450\text{ kg}$ gạo. Ngày thứ hai bán được số gạo gấp đôi ngày thứ nhất. Hỏi sau hai ngày, cửa hàng còn lại bao nhiêu ki-lô-gam gạo?
\shortans{100}
\loigiai{
\textbf{Đáp số:} $100$\\
Ngày thứ hai cửa hàng bán được:
\[
450 \times 2 = 900\text{ (kg gạo)}.
\]
Cả hai ngày cửa hàng đã bán:
\[
450 + 900 = 1\,350\text{ (kg gạo)}.
\]
Số gạo còn lại sau hai ngày là:
\[
1\,450 - 1\,350 = 100\text{ (kg)}.
\]
}
\end{ex}
